\documentclass[11pt,a4paper]{article}

% ---------------------------------------------------------
% PREAMBLE (stable, warning-free, Overleaf-safe)
% ---------------------------------------------------------
\usepackage[utf8]{inputenc}
\usepackage[T1]{fontenc}
\usepackage{lmodern}
\usepackage{amsmath, amssymb, amsfonts}
\usepackage{bm}
\usepackage{geometry}
\usepackage{graphicx}
\usepackage{hyperref}
\usepackage{cite}
\usepackage{authblk}
\usepackage{physics}
\usepackage{mathtools}
\usepackage{textgreek}
\geometry{margin=2.4cm}

\title{\textbf{Companion LXI: Minkowski Functionals of Weak-Lensing Fields in the Relaxation-Driven Cyclic Cosmology}}
\author{Michael Lehmann \\ \texttt{mi.lehmann@gmx.de}}
\date{April 2026}

\begin{document}
\maketitle

\begin{abstract}
Minkowski functionals provide a complete morphological and topological description 
of cosmological fields, capturing information beyond two- and three-point 
statistics. In weak-lensing analyses, they quantify the geometry of the convergence 
field, including area fractions, contour lengths, and Euler characteristics. In the 
standard cosmological model, these functionals are determined by the nonlinear 
matter distribution and the projection of large-scale structure. In the 
Relaxation-Driven Cyclic Cosmology (RDCC), the scale-dependent evolution of the 
gravitational potential modifies the morphology, topology, and connectivity of 
lensing fields in a characteristic way. This Companion develops a continuous 
description of Minkowski functionals in RDCC, deriving how the relaxation scale 
influences the geometry of κ-fields, the topology of excursion sets, and the 
observational signatures in weak-lensing surveys. The resulting framework provides 
a distinctive probe of the RDCC gravitational field beyond correlation functions.
\end{abstract}
% ---------------------------------------------------------
\section{Introduction}

Weak-lensing convergence maps contain rich information about the gravitational 
field. While power spectra and bispectra capture statistical correlations, they do 
not fully describe the geometry or topology of the field. Minkowski functionals 
provide a complete morphological description of excursion sets and are sensitive to 
non-Gaussianity, nonlinear structure formation, and the connectivity of the 
gravitational potential.

In the standard cosmological model, the Minkowski functionals of κ-fields reflect 
the nonlinear matter distribution and the projection of large-scale structure. In 
RDCC, the gravitational potential evolves differently. The relaxation mechanism 
suppresses the decay of small-scale modes and enhances the coherence of large-scale 
modes. This modifies the morphology, topology, and connectivity of κ-fields in a 
scale-dependent manner, producing a distinctive signature in Minkowski functional 
analyses.
% ---------------------------------------------------------
\section{Minkowski Functionals of the Convergence Field}

For a 2D field such as the convergence $\kappa(\hat{\mathbf{n}})$, the Minkowski 
functionals of an excursion set $\kappa > \nu\sigma_{\kappa}$ are:

\begin{align}
    V_{0}(\nu) &= \text{area fraction}, \\
    V_{1}(\nu) &= \text{contour length}, \\
    V_{2}(\nu) &= \text{Euler characteristic}.
\end{align}

These functionals encode:

\begin{itemize}
    \item the abundance of high-κ regions,
    \item the geometry of contours,
    \item the topology of connected structures.
\end{itemize}
% ---------------------------------------------------------
\section{RDCC Modification of the Convergence Field}

In RDCC, the convergence evolves as
\begin{equation}
    \kappa_{\mathrm{RDCC}}(\ell)
    = \kappa(\ell)
      \left[
        1 - \chi_{\kappa}
        \left(\frac{\ell}{\ell_{\mathrm{rel}}}\right)^{\alpha}
      \right].
\end{equation}

The variance becomes
\begin{equation}
    \sigma_{\kappa,\mathrm{RDCC}}^{2}
    = \sigma_{\kappa}^{2}
      \left[
        1 - \chi_{\sigma}
        \left(\frac{\ell}{\ell_{\mathrm{rel}}}\right)^{\alpha}
      \right].
\end{equation}

These modifications alter the morphology of excursion sets.
% ---------------------------------------------------------
\section{Minkowski Functionals in RDCC}

The RDCC Minkowski functionals become
\begin{align}
    V_{0,\mathrm{RDCC}}(\nu)
    &= V_{0}(\nu)
       \left[
         1 - \chi_{0}
         \left(\frac{\ell}{\ell_{\mathrm{rel}}}\right)^{\alpha}
       \right], \\
    V_{1,\mathrm{RDCC}}(\nu)
    &= V_{1}(\nu)
       \left[
         1 - \chi_{1}
         \left(\frac{\ell}{\ell_{\mathrm{rel}}}\right)^{\alpha}
       \right], \\
    V_{2,\mathrm{RDCC}}(\nu)
    &= V_{2}(\nu)
       \left[
         1 - \chi_{2}
         \left(\frac{\ell}{\ell_{\mathrm{rel}}}\right)^{\alpha}
       \right].
\end{align}

These modifications reflect the relaxation-driven changes in morphology and 
topology.
% ---------------------------------------------------------
\section{Topological Signatures of RDCC}

The relaxation mechanism enhances the coherence of large-scale modes. Minkowski 
functionals therefore exhibit:

\begin{itemize}
    \item larger connected high-κ regions,
    \item smoother contours and reduced small-scale curvature,
    \item a shift in the Euler characteristic zero-crossing,
    \item a characteristic turnover near $\ell_{\mathrm{rel}}$,
    \item modified topology of voids and ridges.
\end{itemize}

These signatures provide a direct observational probe of the RDCC gravitational 
field in real-space morphology.
% ---------------------------------------------------------
\section{Conclusion}

Minkowski functionals of weak-lensing fields in the Relaxation-Driven Cyclic 
Cosmology reflect the scale-dependent evolution of the gravitational potential and 
the nonlinear matter distribution. The relaxation mechanism modifies the geometry, 
topology, and connectivity of κ-fields in a scale-dependent manner, producing 
distinctive signatures in weak-lensing surveys. These effects provide a direct 
observational window into the relaxation scale and the temporal structure of the 
RDCC gravitational field. The present Companion establishes the theoretical 
foundation for future studies of topological cosmology, nonlinear structure, and 
the interplay between light and gravity in RDCC.

\bibliographystyle{apsrev4-2}
\nocite{*}
\bibliography{references}


\end{document}
